\begin{abstract}
随着云计算、移动互联网、人工智能和物联网技术的发展，端到端网络传输正在从单一固定终端和局域链路，扩展到跨地域远距离传输、蜂窝/无线/卫星等多接入网络，以及手机、笔记本、台式机和边缘设备等异构终端并存的场景。传统基于丢包的拥塞控制算法（如TCP NewReno、CUBIC）存在拥塞信号滞后、高带宽延迟积（Bandwidth-Delay Product，BDP）网络收敛缓慢、难以区分拥塞程度等局限；基于时延的TCP Vegas和基于带宽/最小RTT模型的BBR则依赖易受路径变化影响的测量量，其与丢包算法共存时的公平性也取决于具体配置。本文提出Swift算法，一种启发式多信号融合自适应网络拥塞控制方法。Swift以确定性规则在确认或拥塞事件上于协议栈之外计算控制决策，不依赖训练样本与神经网络推理，其行为可按规则分支解释，并可在给定配置和随机种子下重复执行。

Swift算法的核心创新收敛为三点：（1）面向远距离与异构终端传输的多信号拥塞状态感知，将ECN状态、拥塞控制状态机状态、拥塞事件类型、RTT和在途字节等信息组织为11维有效观测子空间，并据此区分超时、ECN驱动的窗口缩减与普通丢包；该子空间嵌入15元素状态传输容器，另含4项跨进程路由元数据。（2）启发式反馈自适应的融合控制，根据最小RTT感知的RTT膨胀阈值、性能反馈值相对慢速基线的偏移以及连续增长趋势调节每流乘性增加因子$\alpha\in[0.85,1.30]$，并以滑动时间窗累计确认量和窗口化最大值滤波估计带宽延迟积（Bandwidth-Delay Product，BDP），有界跟踪$\alpha\times BDP$目标窗口。（3）面向不确定拥塞信号的稳定性与安全约束，对普通丢包、ECN和超时分别采用0.70、0.75和0.50的窗口保留因子，并设置最多3次连续缩减、降窗后4个ACK事件冻结、窗口上下界和CA\_LOSS下的陈旧动作作废机制。

本文基于ns-3.40网络仿真平台实现了Swift算法原型，仿真进程与决策进程经ZeroMQ同步请求—应答信道交互（由ns3-gym/OpenGym框架承载15元素状态容器的跨进程传输，仅作为消息通道使用，不含学习逻辑）。实验矩阵覆盖近端高速、共享收敛链路、无线局域网接入、蜂窝移动接入、城域与长途/跨域广域网及LEO/GEO卫星通信共36个场景，并通过\texttt{enable\_udp\_burst}开关构建"纯TCP"与"TCP $+$ 平均负载为瓶颈速率32\%（峰值64\%，50\%占空比）的UDP Burst"两类A/B对照，共288次仿真。全部指标由ns-3 FlowMonitor产物重新导出并\emph{仅在前向数据流上定义}；瓶颈队列统一采用启用ECN标记的RED（$MinTh=0.3Q$、$MaxTh=0.9Q$），四种协议对称开启ECN协商。288次仿真的指标均由ns-3 FlowMonitor前向TCP数据流统计得到，并通过流数完整性、物理取值范围和配置一致性检查。\textbf{稳态实验结果}表明：纯TCP设置下Swift在所考察场景中整体呈现较高的聚合吞吐量，相对逐场景最强基线的增益为$+2.1\%\sim+5.8\%$（均值$+4.1\%$），相关结果覆盖近端高速、无线/蜂窝接入与远距离链路；平均单向时延相对NewReno/CUBIC均值平均低约2.3\%（36个场景中24个更低），相对BBR平均高约0.8\%，个别场景最高约10.4\%；Swift的丢包率与NewReno相当，Jain公平性指数接近1且总体与基线相当。\textbf{跨流量鲁棒性实验}表明：UDP突发下Swift在多数场景中仍呈现吞吐优势（相对最强基线$+2.2\%\sim+5.9\%$），吞吐量保持率均值68.6\%与基线基本一致，未观察到明显吞吐塌缩；GEO、LEO卫星和长途WAN等代表性远距离场景也呈现吞吐优势。上述结果来自每个(场景,协议,设置)的一次确定性运行，不附带跨随机种子置信区间；时延、丢包与公平性存在场景相关的混合结果，不能据此主张全面优于基线。

\keywords{拥塞控制，启发式自适应，多信号融合，显式拥塞通知，自适应参数调优，远距离传输，异构终端}
\end{abstract}

\begin{englishabstract}

End-to-end network transmission now spans long-distance cross-region paths, cellular, wireless, and satellite access networks, and heterogeneous endpoint environments involving phones, laptops, desktops, edge devices, and other terminal classes. Traditional loss-based congestion control algorithms such as TCP NewReno and CUBIC react only after loss, converge slowly in high Bandwidth-Delay Product (BDP) networks, and do not distinguish congestion severity. Delay-based TCP Vegas and model-based BBR rely on delay and bandwidth measurements that can be affected by path variation; their fairness when competing with loss-based algorithms also depends on the network configuration. This paper proposes Swift, a \emph{heuristic} multi-signal adaptive congestion-control algorithm. Swift computes control decisions outside the protocol stack through deterministic rules on acknowledgement or congestion events, without training data or neural-network inference. Its behavior can be inspected branch by branch and repeated under a fixed configuration and random seed.

The core innovations of Swift are consolidated into three aspects: (1) multi-signal congestion-state perception for long-distance and heterogeneous endpoint transmission, using an 11-dimensional effective observation subspace to combine ECN status, Congestion Algorithm (CA) state, congestion events, RTT, and bytes in flight and to classify timeout, ECN-driven reduction, and ordinary loss; the subspace is carried in a 15-element container with four routing metadata fields; (2) heuristic feedback-adaptive fusion control that estimates the Bandwidth-Delay Product (BDP) from cumulative acknowledged bytes over a sliding time window and a windowed max-filter, adjusts per-flow $\alpha\in[0.85,1.30]$ from min-RTT-aware RTT inflation, baseline-relative feedback, and growth trends, and tracks the $\alpha\times BDP$ target with bounded steps; and (3) stability and safety constraints for uncertain congestion signals, including retention factors of 0.70, 0.75, and 0.50 for ordinary loss, ECN, and timeout, at most three consecutive reductions, a four-ACK post-decrease freeze, window bounds, and invalidation of stale pending actions in CA\_LOSS.

We implemented the Swift prototype on the ns-3.40 simulation platform; the simulator and decision processes interact over a synchronous ZeroMQ request--reply channel (the 15-element state container is carried by the ns3-gym/OpenGym framework, used purely as a message channel and containing no learning logic). The experiment matrix spans 36 scenarios covering near-end high-speed links, shared oversubscribed links, wireless LAN access, cellular mobile access, metropolitan/long-haul/inter-region wide area networks, and LEO/GEO satellite links. An \texttt{enable\_udp\_burst} switch constructs paired ``pure TCP'' and ``TCP $+$ On-Off UDP burst'' settings, where the burst has an average offered load of 32\% of the bottleneck rate (64\% peak, 50\% duty cycle), for 288 runs in total. All metrics were re-derived from the ns-3 FlowMonitor artifacts and are defined \emph{on the forward data flows only}; the bottleneck queue is RED with ECN marking ($MinTh=0.3Q$, $MaxTh=0.9Q$), and ECN negotiation is enabled symmetrically for all four protocols. Metrics for the 288 runs were computed from forward TCP flows reported by ns-3 FlowMonitor and checked for flow-count completeness, physical value ranges, and configuration consistency. \textbf{Steady-state results} show that Swift generally provides higher aggregate throughput across the evaluated scenarios, with gains of $+2.1\%$--$+5.8\%$ (mean $+4.1\%$) over the strongest per-scenario baseline. Its mean one-way delay is on average about 2.3\% lower than the NewReno/CUBIC mean but 0.8\% higher than BBR, with increases of up to 10.4\% in individual scenarios. Swift's loss rate is comparable to NewReno, while its Jain fairness index remains close to one and broadly comparable with the baselines. \textbf{Cross-traffic experiments} show a similar throughput pattern in most scenarios under UDP bursts ($+2.2\%$--$+5.9\%$ versus the strongest baseline), while the mean throughput-retention ratio of 68.6\% is comparable with the baselines and no marked throughput collapse is observed. Representative GEO and LEO satellite and long-haul WAN cases also show throughput gains. Each (scenario, protocol, setting) triple corresponds to one deterministic run, so the results have no cross-seed confidence intervals. Delay, loss, and fairness are mixed across scenarios and do not support claims of uniform superiority.

\englishkeywords{Congestion Control, Heuristic Adaptation, Multi-signal Fusion, Explicit Congestion Notification, Adaptive Parameter Tuning, Long-distance Transmission, Heterogeneous Endpoints}

\end{englishabstract}
